\chapter{Spectral Coercivity}

\section{Purpose of Spectral Coercivity}

Spectral coercivity is the analytic form of rigidity used when the relevant
state space is linear, Hilbertian, or operator-theoretic. It says that
movement away from the null space of an operator has positive cost.

In URF language, spectral coercivity supplies a quantitative obstruction:
a defect cannot disappear freely if it lies outside the kernel. The operator
detects the defect, and the spectral gap controls how large the defect must be
relative to its energy.

\section{Operator Framework}

Let \(\mathcal H\) be a Hilbert space with inner product
\[
\langle\cdot,\cdot\rangle.
\]
Let
\[
L:\mathcal D(L)\subseteq \mathcal H\to \mathcal H
\]
be a densely defined self-adjoint nonnegative operator. Nonnegative means
\[
\langle f,Lf\rangle\ge 0
\]
for all \(f\in\mathcal D(L)\).

\begin{definition}[Kernel]
The kernel of \(L\) is
\[
\ker L=\{f\in\mathcal D(L):Lf=0\}.
\]
\end{definition}

\begin{definition}[Orthogonal complement]
The orthogonal complement of the kernel is
\[
(\ker L)^\perp=\{f\in\mathcal H:\langle f,g\rangle=0
\text{ for every }g\in\ker L\}.
\]
\end{definition}

The kernel represents the zero-cost directions. Coercivity concerns the
directions orthogonal to the kernel.

\section{Spectral Gap}

\begin{definition}[Spectral gap]
A nonnegative self-adjoint operator \(L\) has spectral gap \(\lambda>0\) above
its kernel if
\[
\langle f,Lf\rangle\ge \lambda \|f\|^2
\]
for every \(f\in\mathcal D(L)\cap(\ker L)^\perp\).
\end{definition}

The largest admissible \(\lambda\) is the first positive spectral value when
the spectrum is discrete. In finite dimensions, if the eigenvalues are ordered
as
\[
0=\lambda_0\le \lambda_1\le \lambda_2\le\cdots,
\]
then the spectral gap is the first strictly positive eigenvalue.

\section{Coercivity Principle}

\begin{theorem}[Spectral Coercivity]
Let \(L\) be a nonnegative self-adjoint operator with spectral gap
\(\lambda>0\) above its kernel. Then every
\(f\in\mathcal D(L)\cap(\ker L)^\perp\) satisfies
\[
\|f\|^2\le \lambda^{-1}\langle f,Lf\rangle.
\]
\end{theorem}

\begin{proof}
The spectral gap assumption gives
\[
\langle f,Lf\rangle\ge \lambda\|f\|^2.
\]
Since \(\lambda>0\), divide by \(\lambda\) to obtain
\[
\|f\|^2\le \lambda^{-1}\langle f,Lf\rangle.
\]
\end{proof}

This is the basic coercivity estimate. It converts an energy bound into a
norm bound away from the kernel.

\section{Defect Interpretation}

Let \(f\) represent a deformation, error, mode, or defect. If \(f\in\ker L\),
then the operator regards \(f\) as cost-free. If
\(f\perp\ker L\), then spectral coercivity says that \(f\) must pay energy
proportional to its squared norm.

\begin{definition}[Spectral defect]
The spectral defect of \(f\) with respect to \(L\) is
\[
E_L(f)=\langle f,Lf\rangle.
\]
\end{definition}

Under a spectral gap,
\[
E_L(f)\ge \lambda\|f\|^2
\]
for \(f\perp\ker L\). Thus positive-size orthogonal defects have positive
energy cost.

\section{Finite-Dimensional Model}

Let \(L\) be a symmetric positive semidefinite matrix on \(\mathbb R^n\). Let
its eigenvalues be
\[
0=\lambda_0=\cdots=\lambda_{m-1}<\lambda_m\le\cdots\le\lambda_{n-1}.
\]
Then
\[
\ker L=\operatorname{span}\{v_0,\ldots,v_{m-1}\}.
\]
For every \(f\perp\ker L\), write
\[
f=\sum_{j=m}^{n-1}a_jv_j.
\]
Then
\[
\langle f,Lf\rangle
=
\sum_{j=m}^{n-1}\lambda_j a_j^2
\ge
\lambda_m\sum_{j=m}^{n-1}a_j^2
=
\lambda_m\|f\|^2.
\]

\begin{theorem}[Finite-Dimensional Coercivity]
For a symmetric positive semidefinite matrix \(L\), the first positive
eigenvalue \(\lambda_m\) gives the coercivity estimate
\[
\|f\|^2\le \lambda_m^{-1}\langle f,Lf\rangle
\]
for every \(f\perp\ker L\).
\end{theorem}

\begin{proof}
The calculation above gives
\[
\langle f,Lf\rangle\ge\lambda_m\|f\|^2.
\]
Divide by \(\lambda_m>0\).
\end{proof}

\section{Graph Laplacian Example}

Let \(G=(V,E)\) be a finite graph. The graph Laplacian is
\[
L=D-A,
\]
where \(D\) is the degree matrix and \(A\) is the adjacency matrix. For a
function \(f:V\to\mathbb R\),
\[
\langle f,Lf\rangle
=
\sum_{\{u,v\}\in E}(f(u)-f(v))^2
\]
up to the conventional normalization factor.

The kernel consists of functions constant on connected components. If \(G\) is
connected, then \(\ker L\) is the one-dimensional space of constant functions.

\begin{theorem}[Connected Graph Coercivity]
If \(G\) is connected and \(\lambda_1(L)>0\) is the first nonzero Laplacian
eigenvalue, then every mean-zero function \(f\) satisfies
\[
\|f\|^2\le \lambda_1(L)^{-1}\langle f,Lf\rangle.
\]
\end{theorem}

\begin{proof}
For a connected graph, mean-zero functions are orthogonal to the constants,
which form \(\ker L\). Apply spectral coercivity with
\(\lambda=\lambda_1(L)\).
\end{proof}

\section{Cycle Graph}

For the cycle graph \(C_n\), the Laplacian eigenvalues are
\[
\lambda_k=2-2\cos\left(\frac{2\pi k}{n}\right).
\]
The first nonzero eigenvalue is
\[
\lambda_1=2-2\cos\left(\frac{2\pi}{n}\right).
\]
Using
\[
1-\cos x\sim \frac{x^2}{2},
\]
we get
\[
\lambda_1\sim \frac{4\pi^2}{n^2}.
\]

Thus the coercivity constant degenerates as \(n\) grows. Large cycle graphs
have low-cost long-wavelength modes.

\section{Expander Contrast}

A graph family \(G_n\) with Laplacians \(L_n\) has a uniform spectral gap if
there is a constant \(\lambda>0\) such that
\[
\lambda_1(L_n)\ge \lambda
\]
for every \(n\).

\begin{theorem}[Uniform Expander Coercivity]
If \(\lambda_1(L_n)\ge\lambda>0\), then every mean-zero function \(f\) on
\(G_n\) satisfies
\[
\|f\|^2\le \lambda^{-1}\langle f,{L_n}f\rangle.
\]
\end{theorem}

\begin{proof}
Since \(\lambda_1(L_n)\ge\lambda\), the spectral gap estimate gives
\[
\langle f,{L_n}f\rangle\ge \lambda\|f\|^2.
\]
Divide by \(\lambda\).
\end{proof}

This is the analytic sense in which expanders resist low-cost global
deformation.

\section{Connection to Entropy and Capacity}

Spectral coercivity does not by itself imply an EntropyDepth lower bound.
It must be connected to entropy or defect accounting.

A typical bridge has the form
\[
H(x)\le \psi(E_L(f_x))
\]
or
\[
E_L(f_x)\le \psi(H(x)),
\]
depending on the direction needed. The bridge must be stated explicitly.

\begin{definition}[Coercive entropy bridge]
A coercive entropy bridge is an estimate relating an entropy or defect
functional \(H\) to a spectral energy \(E_L\) in a way that allows a capacity
bound on refinement to imply a lower bound on terminal depth.
\end{definition}

Without such a bridge, spectral coercivity remains an analytic estimate rather
than a full URF depth theorem.

\section{Failure Modes}

A spectral-coercivity argument fails if:

\begin{enumerate}
\item the operator is not nonnegative on the claimed domain;
\item the kernel is not identified correctly;
\item the function is not orthogonal to the kernel;
\item the spectral gap vanishes with system size when a uniform gap is needed;
\item no bridge connects spectral energy to entropy or terminal resolution.
\end{enumerate}

Each failure mode is a missing hypothesis, not a contradiction of the
coercivity theorem.

\section{Audit Checklist}

A spectral-coercivity application should provide:

\begin{center}
\begin{tabular}{ll}
\toprule
Field & Required content \\
\midrule
Hilbert space & ambient space and inner product \\
Operator & self-adjoint nonnegative \(L\) \\
Domain & domain \(\mathcal D(L)\) if infinite-dimensional \\
Kernel & zero-cost directions \\
Gap & positive lower bound above the kernel \\
Coercive estimate & norm controlled by energy \\
Bridge & relation to entropy, defect, or terminal resolution \\
Boundary & explicit non-claims \\
\bottomrule
\end{tabular}
\end{center}

\section{Boundary}

This chapter proves standard spectral-coercivity estimates from positive
spectral gaps. It does not prove entropy-depth lower bounds without an
additional entropy bridge, does not prove expansion for any particular graph
family unless the gap is supplied, and does not prove unrestricted Chronos-RR,
unrestricted H4.1/FGL, P vs NP, or any Clay problem.
